% called by main.tex
%
% Norma: máximo 4000 caracteres por idioma, texto plano (sin símbolos),
% estructura clara: introducción/motivación, objetivos, desarrollo y conclusiones.

\section*{Abstract}
\addcontentsline{toc}{section}{Abstract}
\label{sec:abstract}

Choosing a thesis supervisor is one of the most consequential decisions in a student's academic path, yet it is still made with very limited information: word of mouth, lecture impressions, or static topic lists. Paradoxically, the information needed to make this decision well is already public. University research portals describe each faculty member's publications, projects and supervision activity, but in a format designed for one-by-one consultation rather than discovery. This thesis frames supervisor choice as an information retrieval problem and presents \proyecto, a recommendation system that turns the public data of the Scientific Portal of Universidad Politécnica de Madrid into a service that students can query in free text, in any language, to obtain a ranked and explained list of potential supervisors.

The work pursues five objectives: to build an automated and incremental pipeline that acquires and structures the public activity of the \gls{upm} into a graph of researchers, publications, projects and supervised works; to generate an enriched semantic profile per researcher, including normalised topics, expertise domains and an automatically written research biography; to implement and contrast several retrieval strategies, namely sparse lexical, dense semantic, LLM-enriched and hybrid with reciprocal rank fusion; to define an objective and reproducible evaluation framework that requires no manual labelling, using historical supervision records as ground truth, where each supervised project becomes a query whose correct answer is its actual supervisor; and to deliver a usable web prototype, fully deployable on premises with open models, so that no student data ever leaves the institution.

The system was implemented end to end and \pendiente{summarise final corpus figures}. The comparative evaluation over the supervision-based ground truth shows that \pendiente{summarise main quantitative findings: best strategy, contribution of the LLM in the loop, value of hybrid fusion}.

The main conclusions are that institutional public data, processed with locally hosted open language models, suffices to build a measurable and auditable supervisor recommendation service, and that the proposed supervision-based evaluation methodology makes such systems comparable and improvable over time. The pipeline, the evaluation framework and the prototype are released as open source and are directly portable to other institutions with public research portals.

\textbf{Keywords:} recommendation system, expert finding, information retrieval, semantic search, sparse embeddings, dense embeddings, large language models, rank fusion, higher education.

\newpage
\section*{Resumen}
\addcontentsline{toc}{section}{Resumen}
\label{sec:resumen}

La elección del tutor de un trabajo académico es una de las decisiones más determinantes de la trayectoria de un estudiante y, sin embargo, se sigue tomando con información muy limitada: el boca a boca, la impresión de una asignatura o listados estáticos de temas. La paradoja es que la información necesaria para decidir bien ya es pública. Los portales científicos universitarios describen las publicaciones, proyectos y actividad de supervisión de cada investigador, pero en un formato pensado para la consulta ficha a ficha, no para el descubrimiento. Este trabajo plantea la elección de tutor como un problema de recuperación de información y presenta \proyecto, un sistema de recomendación que convierte los datos públicos del Portal Científico de la Universidad Politécnica de Madrid en un servicio que el estudiante consulta en texto libre, en cualquier idioma, para obtener una lista ordenada y justificada de tutores potenciales.

El trabajo persigue cinco objetivos: construir un proceso automatizado e incremental que adquiera y estructure la actividad pública de la \gls{upm} en un grafo de investigadores, publicaciones, proyectos y trabajos supervisados; generar un perfil semántico enriquecido por investigador, con temas normalizados, dominios de especialización y una biografía investigadora redactada automáticamente; implementar y contrastar varias estrategias de recuperación (léxica dispersa, semántica densa, enriquecida con modelos de lenguaje e híbrida con fusión de rankings); definir un marco de evaluación objetivo y reproducible sin etiquetado manual, usando la supervisión histórica como verdad de referencia, donde cada trabajo supervisado se convierte en una consulta cuya respuesta correcta es su supervisor real; y entregar un prototipo web usable, desplegable íntegramente en infraestructura propia con modelos abiertos, de modo que ningún dato salga de la institución.

El sistema se ha implementado de extremo a extremo y \pendiente{resumir las cifras finales del corpus}. La evaluación comparativa sobre la verdad de referencia de supervisión muestra que \pendiente{resumir los hallazgos cuantitativos principales: estrategia ganadora, aportación del LLM y valor de la hibridación}.

Las conclusiones principales son que el dato público institucional, procesado con modelos de lenguaje abiertos alojados localmente, basta para construir un servicio de recomendación de tutores medible y auditable, y que la metodología de evaluación basada en supervisión histórica hace a estos sistemas comparables y mejorables en el tiempo. El proceso de datos, el marco de evaluación y el prototipo se publican como código abierto y son directamente trasladables a otras instituciones con portales científicos públicos.

\textbf{Palabras clave:} sistema de recomendación, búsqueda de expertos, recuperación de información, búsqueda semántica, representaciones dispersas, representaciones densas, modelos de lenguaje de gran tamaño, fusión de rankings, educación superior.
